\section{ADR-0005: Decoupled Subsystem Architecture and Read-Only State Snapshotting}
\label{sec:adr0005}

\begin{adrbox}{ADR-0005: Subsystem Decoupling and State Snapshotting}{Status: \statusbadge{ACCEPTED}}
\textbf{Date:} 2026-11-23 \\
\textbf{Context:} Separation of concerns across gameplay logic, rendering, audio, and platform I/O. \\
\textbf{Deciders:} Core Engine Working Group, Systems Architecture Team \\
\textbf{Status:} Accepted / Production Target
\end{adrbox}

\subsection{Context and Problem Statement}
In monolithic or tightly coupled game engines, visual rendering code and asset I/O are interwoven with core state logic \cite{gregory2018game, sydra2023engine}. This coupling introduces significant vulnerabilities: rendering stalls degrade game physics, graphics API dependencies complicate headless automated testing, and multi-threaded data mutations create data race conditions between render frames and logic ticks \cite{gregory2018game, nystrom2014game}.

The Skip-Logic Game Engine requires complete structural decoupling between mathematical state progression, platform hardware abstraction, rendering views, and asynchronous asset streaming.

\subsection{Decision Drivers}
\begin{enumerate}[label=\bfseries D\arabic*.]
    \item \textbf{Total Decoupling of Logic and View:} The View/Render engine must have zero write access to state progression.
    \item \textbf{Platform Independence:} Isolate OS-specific platform calls behind a uniform Platform Abstraction Layer (PLA).
    \item \textbf{Asynchronous Resource Streaming:} Assets and interactive simulation packages (e.g., PhET physics/math modules) must stream asynchronously without hitching the logic frame.
    \item \textbf{Headless Execution:} Support pure headless CLI execution for CI/CD, mathematical fuzzing, and cloud-based analytics.
\end{enumerate}

\subsection{Decision Outcome}
We establish an immutable **Quad-Layer Architectural Subsystem Separation** paired with an atomic **Read-Only State Snapshotting** protocol.

\begin{table}[htbp]
\centering
\small
\begin{tabularx}{\textwidth}{l p{4.2cm} p{3.8cm} X}
\toprule
\textbf{Engine Layer} & \textbf{Core Responsibilities} & \textbf{Implementation} & \textbf{Primary Data Artifacts} \\
\midrule
\textbf{Logic \& State Engine} & Manages dynamic skip-logic state transitions, Hilbert space updates, and metric integration & Lexlean declarative code compiled via PrismPM to Rust & State Vector $\PsiVec$, Dynamic Operator $\XiOp$, Task DAGs \\
\addlinespace
\textbf{Platform Layer (PLA)} & Low-level OS abstractions, hardware threads, file I/O, windowing, and clock timers & Native Rust cross-platform HAL modules & Platform events, thread pools, raw OS file descriptors \\
\addlinespace
\textbf{View \& Render Engine} & Read-only querying of state frames to drive GPU scene rendering, audio, and UI & Vulkan / wgpu / WebGPU rendering pipelines & Renderable scene graph, audio buffers, UI overlays \\
\addlinespace
\textbf{Resource Subsystem} & Asynchronous streaming, decompression, and caching of asset bundles & Rust async/await background streaming pipeline & Resource handles, cached textures, PhET modules \\
\bottomrule
\end{tabularx}
\caption{Quad-Layer Subsystem Responsibilities and Technologies.}
\label{tab:subsystem_matrix}
\end{table}

\subsubsection{Read-Only State Snapshotting Mechanism}
Communication between the Logic Engine and the View Engine utilizes a triple-buffered, lock-free state publication pattern:
\begin{enumerate}
    \item The Logic Engine writes the new frame state into an internal back buffer $\mathcal{S}_{\text{back}}$.
    \item Upon completing all task graph nodes for frame $N$, the Logic Engine atomically publishes $\mathcal{S}_{\text{back}}$ by swapping an atomic pointer `AtomicPtr<ImmutableStateFrame>`.
    \item The View Engine samples the latest published pointer at its own independent refresh rate (e.g., 144Hz display sync while logic updates at 60Hz), holding a read-only reference without ever stalling the logic pipeline.
\end{enumerate}

\begin{figure}[htbp]
\centering
\begin{tikzpicture}[node distance=2.2cm, auto, >=Stealth]
    % Boxes
    \node [draw, thick, rectangle, fill=primaryblue!10, text width=4.5cm, align=center, rounded corners] (logic) {
        \textbf{Logic \& State Engine}\\
        \footnotesize 60Hz Simulation Tick\\
        Writes $\mathcal{S}_{\text{back}}$ in private arena
    };
    
    \node [draw, ultra thick, dashed, rectangle, fill=accentblue!10, text width=3.5cm, align=center, rounded corners, right=1.5cm of logic] (atomic_pub) {
        \textbf{Atomic Snapshot}\\
        \texttt{AtomicPtr<StateFrame>}\\
        Lock-free swap
    };
    
    \node [draw, thick, rectangle, fill=statusaccepted!10, text width=4.5cm, align=center, rounded corners, right=1.5cm of atomic_pub] (view) {
        \textbf{View / Render Engine}\\
        \footnotesize Variable Refresh (60--144Hz)\\
        Read-Only consumer of $\mathcal{S}_{\text{pub}}$
    };
    
    \node [draw, thick, rectangle, fill=statusproposed!10, text width=11.5cm, align=center, rounded corners, below=1.2cm of atomic_pub] (resource) {
        \textbf{Resource \& Simulation Subsystem (Async PhET Pipeline)}\\
        \footnotesize Background asset streaming $\bullet$ Thread-safe decompression $\bullet$ Zero-copy GPU uploads
    };

    % Arrows
    \draw [->, ultra thick, primaryblue] (logic) -- node[above, font=\footnotesize] {Atomic Publish} (atomic_pub);
    \draw [->, ultra thick, statusaccepted] (atomic_pub) -- node[above, font=\footnotesize] {Read-Only Sample} (view);
    \draw [->, thick, dashed] (resource) -- (logic);
    \draw [->, thick, dashed] (resource) -- (view);
\end{tikzpicture}
\caption{Decoupled Triple-Buffer State Frame Publication Architecture.}
\label{fig:view_decoupling}
\end{figure}

\subsection{Rust Snapshot Publication Interface}

\begin{lstlisting}[language=Rust, caption={Triple-Buffered State Frame Snapshot Publisher}]
use std::sync::atomic::{AtomicPtr, Ordering};

pub struct ImmutableStateFrame {
    pub frame_index: u64,
    pub timestamp_ns: u64,
    pub player_state: PrimeStateVector,
    pub dynamic_operator_norm: f32,
    pub rlg_metric: f32,
    pub efc_metric: f32,
    pub active_sheaf_id: u32,
}

pub struct StateSnapshotPublisher {
    buffers: [Box<ImmutableStateFrame>; 3],
    published_ptr: AtomicPtr<ImmutableStateFrame>,
    write_index: usize,
}

impl StateSnapshotPublisher {
    pub fn publish(&mut self) {
        let ptr = Box::into_raw(self.buffers[self.write_index].clone());
        let old_ptr = self.published_ptr.swap(ptr, Ordering::AcqRel);
        if !old_ptr.is_null() {
            unsafe { drop(Box::from_raw(old_ptr)); }
        }
        self.write_index = (self.write_index + 1) % 3;
    }

    #[inline(always)]
    pub fn read_current(&self) -> &ImmutableStateFrame {
        let ptr = self.published_ptr.load(Ordering::Acquire);
        unsafe { &*ptr }
    }
}
\end{lstlisting}

\subsection{Consequences}
\subsubsection{Positive}
\begin{itemize}
    \item \textbf{Total Visual Isolation:} Heavy render passes or GPU fence waits cannot block state calculation or skip-logic routing.
    \item \textbf{High-Velocity Testing:} Engine logic tests can execute thousands of frames per second in headless mode without initializing graphics devices.
    \item \textbf{Thread Safety Invariants:} Enforced at the compiler level via Rust's immutable referencing guarantees.
\end{itemize}

\subsubsection{Negative and Mitigations}
\begin{itemize}
    \item \textbf{One-Frame Display Latency:} Triple buffering introduces a maximum of one frame display latency ($\sim 16.6\text{ms}$ at 60Hz). \textbf{Mitigation:} Standard game industry trade-off, fully compensated by render-layer extrapolation when required.
\end{itemize}
